% =====================================================================
%  capitulos/cap07.tex — Capítulo 7. De la estructura argumental a las
%  funciones sintácticas
%  Autora: Laura Ferrari
%  Fuente: Capitulo_07_normalizado.docx
% =====================================================================
\capitulo[De la estructura argumental a las funciones sintácticas]%
  {De la estructura argumental a las funciones sintácticas}{Laura Ferrari}

\section{Introducción}

¿Cuáles son las relaciones entre la sintaxis y la semántica? ¿Qué incidencia
posee el léxico en la configuración sintáctica de la oración? ¿Qué tipo de
información está contenida en el nivel léxico? En este capítulo nos proponemos
dar respuesta a estas preguntas. Para ello, partimos del estudio de las
funciones semánticas para ir adentrándonos en las cuestiones sintácticas. En
primer lugar, mostramos cómo el contenido semántico de las palabras que
constituyen una oración determina mayormente sus propiedades combinatorias. Esto
significa que del contenido semántico depende el número de argumentos
seleccionados por un núcleo verbal dentro de su proyección sintagmática. Así, un
verbo como \forma{capturar} necesita de dos SSNN para la buena formación de la
oración que constituye, como en \forma{El policía capturó al ladrón}; en cambio,
\forma{caminar} solo selecciona un SN, como en \forma{Mario camina}. En segunda
instancia, presentamos la \termino{estructura temática}, es decir, las funciones
semánticas de los argumentos que el predicado selecciona, y damos cuenta de la
lista de papeles temáticos. En función de la distinción entre complementos y
adjuntos en el sintagma verbal, desarrollamos las nociones de transitividad
directa e indirecta y distinguimos dos tipos de verbos intransitivos:
inergativos e inacusativos.

\section{Relaciones entre la sintaxis y la semántica}\label{sec:07-sintsem}

\subsection{Predicados y argumentos}\label{sec:07-predicados}

La noción de predicación es semántica y tiene sus orígenes en el análisis
lógico-filosófico. Es un concepto que hace referencia a una condición relativa
al significado de las expresiones. Es decir, la estructura predicativa de las
oraciones es una parte sustancial de su significado. La relación de predicación
se establece entre un predicado y un argumento, como ya se ha señalado en el
capítulo anterior (§\ref{sec:06-predicativa}) respecto del vínculo que se
establece entre el SN sujeto y el SV que predica de él.

\begin{destacado}
Un \termino{predicado} es todo aquel elemento léxico que denota una acción,
proceso o estado y que selecciona un conjunto o red de argumentos en función de
su significado léxico \parencite{rodriguezramalle2005}.
\end{destacado}

Un \termino{argumento} se puede definir como el participante de una acción,
proceso o estado imprescindible para que esta pueda llevarse a cabo.

\begin{destacado}
Los argumentos son «expresiones referenciales que permiten identificar entidades
del mundo extralingüístico» \parencite[102]{ditullio2014}.
\end{destacado}

El ejemplo típico de predicado es el verbo, mientras que los argumentos suelen
ser SSNN.

\begin{exe}
  \ex\label{ex:07-visito} {}[\etq{Arg.~1} El mandatario de Canadá] [\etq{V}
      visitó] [\etq{Arg.~2} el Jardín Japonés].
\end{exe}

\forma{Visitar} en \exref{ex:07-visito} funciona como un predicado que debe
construirse con dos SSNN, que constituyen sus argumentos. El evento de
\forma{visitar} relaciona dos argumentos: x (= \forma{el mandatario de Canadá})
e y (= \forma{el Jardín Japonés}). Para que una oración esté completa, el
elemento seleccionador, en este caso el verbo, debe haberse combinado con todos
los argumentos que selecciona.

\begin{ejercicio}\label{ej:07-agramatical}
Explique la agramaticalidad o el carácter ambiguo de las siguientes oraciones,
según se muestra en los ejemplos resueltos.

\begin{exe}
  \ex María prestó a su hijo.
  \ex La comisión respaldó durante el receso parlamentario.
  \ex José colocó.
  \ex El viento provocó.
  \ex El vocero presidencial anunció gelatina.
  \ex El presidente reclamó a los bancos.
  \ex Los cóndores habitan.
  \ex Lucía va todas las semanas.
\end{exe}

\begin{resuelto}[Ejemplos resueltos]
1. Es ambigua: a) se puede interpretar que lo prestado por \forma{María} es
\forma{su hijo} y queda implícito \sig{a quién}; b) es agramatical porque falta
el argumento referido a \sig{lo que María le prestó a su hijo}. 2. Es
agramatical porque falta el argumento referido a \sig{lo respaldado por la
comisión}.
\end{resuelto}
\end{ejercicio}

\begin{ejercicio}\label{ej:07-seleccion}
Algunas de las oraciones que siguen, por diversas razones, no respetan los
requisitos de selección de los verbos que las conforman. Haga una lectura atenta
y responda lo pedido en las consignas, según se muestra en los ejemplos
resueltos de abajo.

\begin{consignas}
  \item Explique el porqué de la buena o mala formación de cada una de las
        oraciones.
  \item Luego especifique la estructura argumental de \forma{regresar} y
        \forma{llegar}.
\end{consignas}

\begin{exe}
  \ex El viajero regresó contento a su casa y trajo muchos regalos a sus
      familiares.
  \ex El viajero llegó contento a su casa y trajo muchos regalos.
  \ex Llegó contento a su casa y trajo muchos regalos.
  \ex El viajero llegó oblongo a su casa y trajo muchos regalos a sus familiares.
  \ex El viajero llegó contento su casa y trajo muchos regalos a sus familiares.
  \ex El viajero llegó contento a su aceituna y trajo muchos regalos a sus
      familiares.
\end{exe}

\begin{resuelto}[Ejemplos resueltos]
a) 1 es gramatical porque ambos verbos tienen todos sus argumentos. 4 es
agramatical porque \forma{oblongo} no es compatible desde el punto de vista
léxico-semántico con el sujeto ni con el verbo \forma{llegar}. b)
\forma{Llegar} [Arg.~1 SN, Arg.~2 SP].
\end{resuelto}
\end{ejercicio}

\subsubsection{La valencia verbal}\label{sec:07-valencia}

El número de argumentos que necesita un verbo para que esté completo o saturado
se denomina \termino{valencia verbal}. Pero ¿cuál es el número de argumentos de
un predicado? El número y el tipo de argumentos de un predicado dependen de cómo
se conceptualiza lingüísticamente el evento. Y ¿cuántos argumentos puede exigir
un verbo? Para \textcite{hale1998}:

\begin{cita}
Los verbos de una lengua natural, generalmente la más rica categoría al
respecto, son extremadamente limitados en la variedad y complejidad de las
estructuras argumentales que despliegan, y estas se conforman a una tipología
altamente restricta. \textbf{Pocos verbos tienen más de tres argumentos}\ldots{}
(Traducción nuestra)
\end{cita}

Según el número de argumentos que exigen, los verbos se clasifican en:

\begin{itemize}
  \item \termino{cero-ádicos} o \termino{avalentes}: \forma{Continúa granizando
        (en el sur)};\footnote{Los paréntesis indican los adjuntos.}
        \forma{Llovió (torrencialmente) (toda la noche)};
  \item \termino{monádicos} o \termino{monovalentes}: \forma{[El ministro de
        economía] viajó (ayer)}; \forma{Hay [café]};
  \item \termino{diádicos} o \termino{bivalentes}: \forma{[El juez] redactó [la
        sentencia]}; \forma{[Mi vecino] había perdido [las llaves]};
  \item \termino{triádicos} o \termino{trivalentes}: \forma{[El fiscal] entregó
        [sus informes] [a la justicia]}; \forma{[El periodista] solicitó
        [información] [a la embajada]}.
\end{itemize}

Los predicados seleccionan sus argumentos en función de su significado. El
conjunto de argumentos de un predicado representa un esqueleto de su propia
definición y constituye su \termino{estructura argumental}.

\begin{ejercicio}\label{ej:07-valencia}
Lea atentamente el siguiente fragmento de «Emma Zunz», de J.~L.~Borges, y luego
resuelva lo pedido en las consignas, según se muestra en el ejemplo resuelto de
abajo.

% Fragmento literario protegido: se incorpora a mano desde el original.
\textoaincorporar{el primer párrafo de «Emma Zunz», de J.~L.~Borges (\textit{El
Aleph}), desde «El catorce de enero de 1922» hasta «\ldots{}que se dirigía a la
hija del muerto»}

\begin{consignas}
  \item Según la cantidad de argumentos que selecciona, determine a qué clase
        pertenece cada verbo (monádico, triádico, etc.).
  \item Luego identifique los argumentos seleccionados por cada verbo.
\end{consignas}

\begin{resuelto}
\forma{Hallar}: a) es un V triádico porque selecciona tres argumentos; b) los
argumentos seleccionados son \forma{Emma Zunz} (sujeto), \forma{una carta,
fechada en el Brasil} (COD) y \forma{en el fondo del zaguán} (CCL).
\end{resuelto}
\end{ejercicio}

\subsection{Complementos y adjuntos}\label{sec:07-compadj}

Como ya se introdujo en el \capref{6}, §\ref{sec:06-externo}, existe una
diferencia entre el \termino{argumento externo} o sujeto, que por su relación
estructural prominente es el que mantiene con el verbo la relación de
predicación, y los argumentos internos o complementos verbales, que permanecen
dentro del SV \parencite{williams1980}.\pendiente{El original dice «Williams
1981»}

Los argumentos internos al SV seleccionados por el predicado se denominan
\termino{complementos}. Los complementos y el argumento externo (sujeto)
introducen información exigida por el significado de los predicados, como
\forma{el señor Maier} (argumento externo / sujeto) y \forma{una fuerte dosis de
veronal} (argumento interno / COD) en \exref{ex:07-maier-a}. Por otra parte, los
modificadores no seleccionados por el predicado se denominan
\termino{adjuntos}, como los que aparecen entre paréntesis en
\exref{ex:07-maier-a}. Generalmente los adjuntos pueden omitirse y la oración
sigue siendo gramatical \exref{ex:07-maier-b}. Los adjuntos también pueden
agregarse en una oración con gran libertad \exref{ex:07-maier-c}; en cambio,
cada complemento puede ser realizado exclusivamente por un único argumento. En
otras palabras, un verbo puede tener varios complementos, siempre y cuando
realicen diferentes funciones sintácticas. Como se puede apreciar,
\exref{ex:07-maier-d} es agramatical porque hay dos SSNN no coordinados que
realizan la función de OD. No obstante, algunos argumentos pueden quedar
implícitos, esto quiere decir que forman parte de la estructura argumental del
verbo, pero no se manifiestan necesariamente en la sintaxis
\exref{ex:07-fruta}.

\begin{exe}
  \ex\label{ex:07-maier}
    \begin{xlist}
      \ex\label{ex:07-maier-a} {}[El señor Maier] había ingerido (por error) [una
          fuerte dosis de veronal] y había fallecido (el tres del corriente) (en
          el hospital de Bagé).
      \ex\label{ex:07-maier-b} {}[El señor Maier] había ingerido [una fuerte
          dosis de veronal].
      \ex\label{ex:07-maier-c} {}[El señor Maier] había ingerido (por error) [una
          fuerte dosis de veronal] (con gran repugnancia) (esa mañana).
      \ex[*]{\label{ex:07-maier-d} [El señor Maier] había ingerido [una fuerte
          dosis de veronal] [veneno para ratas].}
    \end{xlist}
  \ex\label{ex:07-fruta} ---¿Comés fruta?\\
      ---No, no como. (El COD de \forma{comer} está implícito).
\end{exe}

\begin{ejercicio}\label{ej:07-compadj}
Lea atentamente el siguiente fragmento de «Emma Zunz» y resuelva lo pedido en
las consignas, según se muestra en el modelo de abajo.

% Fragmento literario protegido: se incorpora a mano desde el original.
\textoaincorporar{el fragmento de «Emma Zunz» que va de «En abril cumpliría
diecinueve años» hasta «\ldots{}se acostó y se obligó a dormir»}

\begin{consignas}
  \item Identifique los complementos internos al sintagma verbal y los adjuntos
        presentes en los predicados.
  \item Indique si alguno de los verbos posee un argumento implícito.
\end{consignas}

\begin{resuelto}
a) \forma{Cumpliría}: \forma{diecinueve años} (complemento) y \forma{en abril}
(adjunto). \forma{Inspiraban}: \forma{los hombres}, \forma{le} y \forma{un temor
casi patológico} (complementos); \forma{aún} (adjunto). b) \forma{Cumpliría}
tiene el argumento sujeto tácito o implícito.
\end{resuelto}
\end{ejercicio}

\section{La estructura temática}\label{sec:07-tematica}

Además de la información acerca de la cantidad y características de los
argumentos requeridos por un verbo, es necesario reconocer que los distintos
argumentos cumplen diferentes funciones semánticas, conocidas como papeles
temáticos, es decir, que los predicados asignan un valor semántico a los
argumentos que seleccionan. ¿Cómo definimos los papeles temáticos?

\begin{destacado}
Los \termino{papeles temáticos} dan cuenta de la relación establecida por cada
participante en el evento descripto por el verbo.
\end{destacado}

\subsection{La estructura temática como esqueleto del significado: los papeles temáticos}\label{sec:07-papeles}

La \termino{estructura temática} de un predicado es la variante enriquecida de
su estructura argumental, es decir, contiene información sobre la semántica de
los argumentos que el predicado selecciona.

Los predicados otorgan a sus argumentos distintos papeles temáticos que expresan
su participación en el evento. Si bien no hay acuerdo generalizado respecto de
cuántos y cuáles son los papeles temáticos, en la siguiente lista incluimos los
más reconocidos y en los que hay mayormente consenso.

\begin{itemize}
  \item \textbf{Agente}: el que voluntariamente causa o realiza una acción, como
        \forma{Emma} en \forma{Emma recogió el papel}. El agente es típicamente
        \sig{animado}. Cuando el causante de la acción no tiene el rasgo
        \sig{animado} se identifica como \textbf{Causa} o \textbf{Fuerza}, como
        \forma{el fuego} en \forma{El fuego destruyó la cabaña}, o como
        \textbf{Instrumento} ---el medio u objeto empleado para realizar la
        acción---, como \forma{el revólver} en \forma{Emma mató a Loewenthal con
        un revólver}.
  \item \textbf{Tema}: el objeto efectuado, afectado o trasladado, o cuya
        existencia o posición se predica, como \forma{la verja} en \forma{Emma
        empujó la verja}. Cuando es \sig{animado} suele denominarse
        \textbf{Paciente}, como \forma{el marinero} en \forma{Emma buscó al
        marinero}.
  \item \textbf{Experimentante}: el que experimenta un estado psicológico o
        cognitivo determinado, como \forma{Emma} en \forma{Emma sintió asco}.
  \item \textbf{Beneficiario}: el que recibe beneficio o daño
        ---\textbf{Maleficiario}--- de un suceso, como \forma{Perla} en
        \forma{Perla perdió su trabajo}.
  \item \textbf{Locativo}: el que designa una ubicación en relación con el
        evento, como \forma{en Almagro} en \forma{Pedro vivía en Almagro}.
  \item \textbf{Fuente} (u \textbf{Origen}): entidad o lugar a partir del cual se
        origina la actividad, como \forma{su socio} en \forma{Recibió el arma de
        su socio} o \forma{el Río de la Plata} en \forma{El viaje empieza en el
        Río de la Plata}.
  \item \textbf{Meta}: la entidad o lugar hacia el que se dirige la actividad,
        como \forma{la empleada} en \forma{Entregó los informes a la empleada} o
        \forma{la estación} en \forma{La avenida llegaba hasta la estación}.
  \item \textbf{Vía}: el lugar de tránsito del movimiento, como \forma{por la
        avenida 9 de Julio} en \forma{Llegó hasta el puerto por la avenida 9 de
        Julio}.
  \item \textbf{Cantidad} o \textbf{Medida}: la medida espacial o temporal del
        evento, como \forma{veinte centímetros} en \forma{Creció veinte
        centímetros}.
\end{itemize}

\begin{destacado}
\textbf{Criterios de reconocimiento}
\begin{itemize}
  \item Cada argumento recibe un solo papel temático.
  \item Cada papel temático se asigna a un solo argumento.
  \item Los papeles temáticos no coinciden con las funciones sintácticas
        tradicionales.
  \item Los adjuntos no forman parte de la estructura argumental.
\end{itemize}
\end{destacado}

\begin{ejercicio}\label{ej:07-papeles}
En el fragmento de «Emma Zunz» del ejercicio~\ref{ej:07-valencia}, señale los
papeles temáticos de los sintagmas subrayados e indique cuál es el predicado que
los selecciona, según se muestra en el ejemplo resuelto de abajo.

\textoaincorporar{los subrayados del primer párrafo de «Emma Zunz»: \textit{de la
fábrica de tejidos Tarbuch y Loewenthal; una carta; su padre; el sello y el
sobre; la letra desconocida; Emma; el señor Maier; un compañero de pensión de su
padre; la noticia; a la hija del muerto}}

\begin{resuelto}
\forma{De la fábrica de tejidos Tarbuch y Loewenthal} tiene el papel temático de
fuente y es seleccionado por el predicado \forma{volver}.
\end{resuelto}
\end{ejercicio}

\begin{recuadro}[Para ampliar]
\textcite{bosque2009} señalan algunos beneficios de trabajar con estructuras
temáticas, pero también advierten que este enfoque presenta algunos
inconvenientes.

\medskip\textbf{Ventajas de postular estructuras temáticas}
\begin{itemize}
  \item La estructura temática muestra la correspondencia entre papeles
        temáticos y construcciones estructurales canónicas; por ejemplo, SN para
        Agente y para Tema (encabezado por \forma{a}) y SP para Meta.
  \item La estructura temática permite deslindar la información sintáctica y la
        información semántica.
  \item Las estructuras temáticas permiten considerar clases semánticas
        naturales. Por ejemplo, podemos analizar los verbos del tipo [A, T, L]
        (Agente, Tema, Locativo) como una clase que agrupa a verbos tales como
        \forma{situar, colocar, ocultar, poner}, en lugar de asignarle
        individualmente una estructura temática a cada predicado.
  \item Las estructuras temáticas no informan sobre los adjuntos, sino sobre los
        argumentos. Algunas nociones pueden manifestarse como complementos con
        ciertos verbos y como adjuntos con otros predicados, lo que confirma que
        tiene importancia caracterizar a estos modificadores en relación con el
        predicado en el que participan. Por ejemplo, si comparamos
        \forma{Colocó los libros en la mesa del comedor} con \forma{Hizo la tarea
        en la mesa del comedor}, \forma{en la mesa del comedor} es complemento
        (argumento) en el primer caso y adjunto en el segundo.
\end{itemize}

\textbf{Inconvenientes de postular estructuras temáticas}
\begin{itemize}
  \item ¿Cuántos papeles temáticos tenemos exactamente? Lo que se pone en duda
        actualmente es que cada argumento de cada predicado tenga que recibir
        una etiqueta semántica que represente un concepto definido entre las
        nociones que articulan la teoría de la gramática.
  \item Algunos papeles temáticos son demasiado abarcadores. El ejemplo más
        claro lo representa el papel de Tema, que refiere tanto a OD canónicos,
        como en \forma{Comió una milanesa}, como a nociones menos específicas,
        como en \forma{Miró la pantalla de su celular}, donde el objeto no es ni
        afectado, ni efectuado, ni trasladado.
  \item ¿Cómo estar seguros de que identificamos el papel temático correcto? No
        siempre hay garantías demasiado claras y existen casos que pueden ser
        menos transparentes para su reconocimiento.
\end{itemize}
\end{recuadro}

En la estructura temática, los dos papeles de mayor centralidad son el agente y
el tema, que son los pilares de la oración transitiva, si bien también están
presentes en otros tipos de esquemas oracionales que revisaremos.

\subsection{La noción de agentividad}\label{sec:07-agentividad}

Las propiedades que permiten caracterizar a los agentes prototípicos son
\sig{intencionalidad}, \sig{control consciente sobre la acción}, \sig{origen,
causa o iniciador del proceso} y \sig{autonomía} (es decir, que el agente
preexiste al evento). Por otra parte, para identificar el agente se puede
aplicar una paráfrasis con el verbo \forma{hacer}. Así, en \exref{ex:07-juan},
esta prueba confirma que \forma{Juan} es agente.

\begin{exe}
  \ex\label{ex:07-juan}
    \begin{xlist}
      \ex {}[Juan] trabajó toda la tarde.
      \ex Lo que hizo Juan fue trabajar toda la tarde.
    \end{xlist}
\end{exe}

En cambio, en \exref{ex:07-telon}, si aplicamos esta prueba comprobamos que
\forma{el telón} no es agente, sino tema.

\begin{exe}
  \ex\label{ex:07-telon}
    \begin{xlist}
      \ex Cayó [el telón].
      \ex[*]{Lo que hizo el telón fue caer.}
    \end{xlist}
\end{exe}

Asimismo, solo una oración que posee un sujeto agentivo admite las paráfrasis
con construcciones finales, que indican la intencionalidad de quien realiza la
acción.

\begin{exe}
  \ex {}[Juan] trabajó para aumentar sus ingresos.
  \ex {}[Juan] camina para hacer ejercicio.
  \ex[*]{[El telón] cayó para terminar la función.}
  \ex[*]{[Juan] palideció para hacer ejercicio.}
\end{exe}

Los sujetos de los verbos de acción normalmente representan una «causa
externa», es decir que la transformación que puede sufrir el tema o paciente
tiene su origen en una causa que ejerce control sobre el evento y que puede ser
un agente que realiza deliberadamente esa acción. Por ejemplo, en \forma{Los
chicos arreglaron su cuarto}, el sujeto agente ejerce el control sobre el evento
y realiza deliberadamente la acción de \forma{arreglar el cuarto}. También puede
darse el caso de una causa desencadenante no animada que provoque una situación
determinada, como en \forma{El viento derribó el toldo}. O incluso un
instrumento que, por medio de un agente no manifiesto, origina un determinado
estado de cosas, como en \forma{El camión embistió el portón de hierro}, que
supone un conductor implícito.

\begin{ejercicio}\label{ej:07-agente}
Indique en el texto siguiente si los SSNN subrayados poseen el papel temático de
agente o no. Justifique su respuesta a partir de la aplicación de las pruebas
pertinentes, según se muestra en el ejemplo resuelto de abajo.

% Texto periodístico protegido: se incorpora a mano desde el original.
\textoaincorporar{el texto sobre la conferencia de David Hilbert en la Sorbona
(1900) y «los problemas de Hilbert», con los SSNN subrayados: \textit{un alemán
llamado David Hilbert; Hilbert; la legendaria lista}}

\begin{resuelto}
\forma{Un alemán llamado David Hilbert} es agente. Acepta la paráfrasis con
\forma{hacer} (\forma{Lo que hizo X fue ponerles a los asistentes la tarea de
matemáticas\ldots}) y la construcción final (\forma{Un alemán llamado David
Hilbert les puso a los asistentes la tarea de matemáticas (\ldots) para lograr
su propósito}).
\end{resuelto}
\end{ejercicio}

\section{El esquema transitivo}\label{sec:07-transitivo}

La noción de \termino{objeto directo} o \termino{complemento objeto directo}
(COD) es una herencia de la tradición gramatical, para la cual «el complemento
directo recibe la acción del verbo, e incluso [\ldots] esta pasa al objeto o se
manifiesta en él» \parencite[§\,34.1e]{rae2025}. Por ejemplo, en
\exref{ex:07-albaniles}, el SN \forma{la pared} constituye el objeto efectuado
por la acción expresada por el verbo y cumple la función sintáctica de COD.

\begin{exe}
  \ex\label{ex:07-albaniles} Los albañiles construyeron [la pared].
\end{exe}

Desde el punto de vista categorial, el objeto directo es un SN (\forma{El
cocinero preparó \textbf{la comida}}) o una oración
subordinada\footnote{La noción de oración subordinada se trata en el
\capref{13}.} (\forma{Los directores dijeron \textbf{que no les había gustado el
reglamento}}).

La definición tradicional de transitividad se basa en la existencia de acciones
realizadas por un agente (sujeto) que pasan o se trasladan a un objeto tema o
paciente. Sin embargo, existe una importante vaguedad e imprecisión en estos
conceptos, dado que no son aplicables a todas las estructuras transitivas
posibles. Por ejemplo, en una oración como \forma{María vio un avión}, el sujeto
posee el papel temático de experimentante y no ejerce ninguna acción que pueda
trasladarse al objeto.

Definimos, entonces, un \termino{verbo transitivo} como aquel que selecciona su
complemento y lo identifica con la marca de caso acusativo. Ese complemento
cumple la función sintáctica de COD. Léxicamente, los verbos transitivos
requieren la presencia de un complemento que complete su significado.

Desde el punto de vista sintagmático, la relación estrecha entre el V y el COD
puede observarse en el diagrama arbóreo \exref{ex:07-arbol-cod}.

\begin{exe}
  \ex\label{ex:07-arbol-cod}
      \begin{forest} gramatica
        [SV
          [SN\\(Suj)\\\textit{los albañiles}]
          [V\prima
            [V\\\textit{construyeron}]
            [SN\\(COD)\\\textit{la pared}]]]
      \end{forest}
\end{exe}

\subsection{Objeto efectuado y objeto afectado}

Existe otra distinción gramaticalmente relevante entre los casos en que el
objeto es \termino{efectuado} por la acción expresada por el verbo y aquellos en
que solamente es \termino{afectado} por ella, es decir, los casos en que la
acción se aplica a un objeto ya existente. Veamos la diferencia con un ejemplo
como \exref{ex:07-mesada}.

\begin{exe}
  \ex\label{ex:07-mesada}
    \begin{xlist}
      \ex\label{ex:07-mesada-a} Rubén construyó [la mesada].
      \ex\label{ex:07-mesada-b} Rubén pintó [la mesada].
    \end{xlist}
\end{exe}

En \exref{ex:07-mesada-a}, \forma{la mesada} refiere a un objeto que se
materializa a través de la acción del verbo \forma{construir}; en
\exref{ex:07-mesada-b}, en cambio, \forma{la mesada} señala un objeto
preexistente, que se modifica por la acción indicada por el verbo.

Los objetos directos admiten la presencia de un clítico acusativo (\forma{lo,
la, los, las}) que sustituye o dobla el objeto directo, como se muestra en la
paráfrasis siguiente:

\begin{exe}
  \ex Rubén [la] construyó.
\end{exe}

\subsection{El objeto directo animado y el marcado diferencial de objeto (MDO)}\label{sec:07-mdo}

En español, algunos objetos directos poseen una marca especial \forma{a}. Veamos
algunos ejemplos:

\begin{exe}
  \ex
    \begin{xlist}
      \ex Los docentes recibieron [a los alumnos] en el primer día de clase.
          (OD \sig{animado})
      \ex[*]{Los docentes recibieron [los alumnos] en el primer día de clase.}
    \end{xlist}
  \ex
    \begin{xlist}
      \ex Los empleados recibieron [la intimación]. (OD \sig{inanimado})
      \ex[*]{Los empleados recibieron [a la intimación].}
    \end{xlist}
  \ex\label{ex:07-perro}
    \begin{xlist}
      \ex Mariano encontró [un perro abandonado]. (OD \sig{indeterminado})
      \ex[?]{\label{ex:07-perro-b} Mariano encontró [a un perro abandonado].}
      \ex Mariano encontró [al perro abandonado].
    \end{xlist}
\end{exe}

En términos generales, el objeto directo lleva la marca de función \forma{a}
cuando refiere a una entidad animada, determinada y específica, aunque puede
haber excepciones, como \exref{ex:07-perro-b}. ¿Por qué \exref{ex:07-perro-b}
nos resulta extraña o poco aceptable? La oración resulta anómala porque lleva la
preposición \forma{a} un objeto directo inespecífico: \forma{un perro
abandonado}.

Algunos usos impersonales de verbos tales como \forma{haber} (\forma{No hay
plata}) y \forma{hacer} (\forma{Hace más de treinta años}) seleccionan un
complemento objeto directo de carácter marginal, ya que la sustitución por el
clítico acusativo (\forma{lo, la}) es solo posible en ciertos contextos
marcados.

\begin{ejercicio}\label{ej:07-mdo}
Señale en los textos siguientes las características de los objetos directos
subrayados (± animado, ± específico). Fundamente su respuesta, según se muestra
en los ejemplos resueltos de abajo.

% Textos protegidos: se incorporan a mano desde el original.
\textoaincorporar{Texto 1: el fragmento sobre Cantor y el infinito, con los COD
subrayados \textit{el significado del infinito; precisión matemática; un solo
infinito}. Texto 2: el fragmento de «Emma Zunz» que empieza «Recordó veraneos en
una chacra», con los COD subrayados \textit{veraneos; a su madre; la casita de
Lanús; los amarillos losanges de una ventana; el auto de prisión; el oprobio; los
anónimos}}

\begin{resuelto}[Ejemplos resueltos]
\forma{El significado del infinito}: COD \sig{inanimado} y \sig{específico}.
\forma{Un solo infinito}: COD \sig{inanimado} e \sig{inespecífico}, seleccionado
por un verbo impersonal de significado \sig{existencial}.
\end{resuelto}
\end{ejercicio}

\subsection{Los complementos preposicionales}\label{sec:07-regimen}

Si bien las gramáticas tradicionales solían incluir los complementos
preposicionales entre los circunstanciales, no todos ellos se comportan en
forma análoga. En el caso del \termino{complemento de régimen}, se trata de un
modificador seleccionado por el núcleo del predicado, es decir, tiene carácter
argumental. Por otra parte, el núcleo del predicado selecciona también la
preposición que lo encabeza. Veamos algunos ejemplos:

\begin{exe}
  \ex\label{ex:07-walmart} La marca Walmart se despide [del mercado argentino].
  \ex\label{ex:07-discutio} El ministro discutió [con los empresarios].
  \ex\label{ex:07-penso} Juan pensaba [en su madre].
  \ex\label{ex:07-quejo} Juana se quejó [de su madre].
\end{exe}

El complemento de régimen preposicional no puede ser sustituido ni duplicado por
un pronombre clítico acusativo o dativo. Algunos verbos, como \forma{despedir},
\forma{pensar} y \forma{discutir}, se comportan como transitivos en oraciones
tales como \forma{María despidió a su padre}, \forma{El ministro discutió sus
argumentos} o \forma{Juan pensaba la solución del problema}, en las que
seleccionan un COD, a diferencia de oraciones como \exref{ex:07-walmart},
\exref{ex:07-discutio} y \exref{ex:07-penso}, en las que se construyen con un
complemento de régimen. En cambio, el verbo \forma{quejar(se)}, en
\exref{ex:07-quejo}, no posee esa alternancia y siempre se construye con un
complemento encabezado por la preposición \forma{de}, aunque en ciertos
contextos el complemento puede quedar implícito: \forma{Juana se quejó todo el
tiempo}.

\subsection{Otros complementos}

Algunos verbos seleccionan complementos encabezados por una preposición cuyo
significado determina el papel temático del complemento que lo acompaña. Es el
caso de los \termino{complementos locativos} o \termino{de modo}:

\begin{exe}
  \ex\label{ex:07-volvia} Volvía [de Mar del Plata] por aquel entonces.
  \ex\label{ex:07-conduce} Conduce [a su manera].
\end{exe}

Los complementos señalados en \exref{ex:07-volvia} y \exref{ex:07-conduce} están
seleccionados por los verbos \forma{volver} y \forma{conducir}, respectivamente.
Se trata de complementos circunstanciales y no debemos confundirlos con los
adjuntos circunstanciales, que no son requeridos por el verbo, como, por
ejemplo, en:

\begin{exe}
  \ex\label{ex:07-trabaja} Trabaja (en Mar del Plata) (durante el verano).
  \ex\label{ex:07-pinto} Pintó la silla (a su manera).
\end{exe}

En \exref{ex:07-trabaja} y \exref{ex:07-pinto}, las construcciones entre
paréntesis funcionan como adjuntos circunstanciales de lugar y modo,
respectivamente.

\subsection{Esquemas ditransitivos}\label{sec:07-ditransitivos}

Son aquellos esquemas oracionales en los que el verbo selecciona dos argumentos
internos: uno directo y otro indirecto. En una oración como \forma{Los
socorristas pidieron [ayuda] [a los bomberos]}, los sintagmas entre corchetes
constituyen los dos complementos de objeto: el primero, el directo, y el
segundo, el indirecto. Típicamente, existen tres subclases de verbos que
constituyen esquemas ditransitivos: los verbos de comunicación
\exref{ex:07-anuncio}, los de transferencia \exref{ex:07-presto} y los de
influencia \exref{ex:07-solicito}. Veamos los ejemplos siguientes:

\begin{exe}
  \ex\label{ex:07-anuncio} El directorio anunció las medidas [a sus miembros].
  \ex\label{ex:07-presto} Prestó su bolsa de dormir [a su sobrino].
  \ex\label{ex:07-solicito} Solicitó [a sus jefes] una licencia.
\end{exe}

Los objetos indirectos admiten la presencia de un clítico dativo (\forma{le} o
\forma{les}), que sustituye o dobla el objeto indirecto, como se muestra en las
paráfrasis siguientes:

\begin{exe}
  \ex {}[Les] anunció las medidas.
  \ex {}[Le] prestó su bolsa de dormir.
  \ex {}[Les] solicitó una licencia.
\end{exe}

En el diagrama que presentamos a continuación se muestra la configuración
estructural del complemento objeto indirecto (COI) en el sintagma verbal y la
diferencia con la configuración estructural del COD.

\begin{exe}
  \ex\label{ex:07-arbol-coi} El directorio anunció las medidas a sus miembros.\\
      \begin{forest} gramatica
        [SV
          [V\prima
            [V\prima
              [V\\\textit{anunció}]
              [SN\\(COD)\\\textit{las medidas}]]
            [SP\\(COI)\\\textit{a sus miembros}]]]
      \end{forest}
\end{exe}

Otra clase de verbos ditransitivos (diádicos o bivalentes) son los llamados
verbos de «colocación», que seleccionan dos argumentos internos: un tema (COD) y
un locativo (CCL), como se expresa en los ejemplos que presentamos a
continuación:

\begin{exe}
  \ex Colocó estantes en la cocina.
  \ex Colgó los cuadros en la pared del dormitorio.
\end{exe}

\begin{ejercicio}\label{ej:07-seleccion2}
Lea atentamente el siguiente fragmento de «Emma Zunz» y responda lo pedido según
se muestra en los ejemplos resueltos de abajo.

\begin{consignas}
  \item Señale qué verbos seleccionan los sintagmas y pronombres clíticos
        subrayados.
  \item Relacione los sintagmas y pronombres clíticos subrayados con los verbos
        que los seleccionan.
\end{consignas}

% Fragmento literario protegido: se incorpora a mano desde el original.
\textoaincorporar{los dos fragmentos de «Emma Zunz» que van de «Emma, desde
1916, guardaba el secreto» hasta «\ldots{}el sabor de la victoria y de la
justicia» y de «Dio al fin con hombres del Nordstjärnan» hasta «\ldots{}a una
puerta que se cerró», con los subrayados: \textit{a nadie; a su mejor amiga,
Elsa Urstein; le} (dos veces)\textit{; hombres del Nordstjärnan; por otro; a una
puerta; a un turbio zaguán; a una escalera tortuosa; a un vestíbulo; a un
pasillo; a una puerta que se cerró}}

\begin{resuelto}[Ejemplos resueltos]
\forma{A nadie} y \forma{a su mejor amiga, Elsa Urstein} son seleccionados por
\forma{revelar}. a) El primer \forma{le} es seleccionado por \forma{inspirar}.
El segundo \forma{le} es seleccionado por el verbo \forma{deparar}. b)
\forma{Revelar} es un verbo de habla y selecciona un complemento objeto
indirecto, y \forma{deparar} es un verbo psicológico y selecciona un complemento
objeto indirecto.
\end{resuelto}
\end{ejercicio}

\subsection{Esquemas intransitivos: inergativos e inacusativos}\label{sec:07-intransitivos}

Los verbos intransitivos se han caracterizado tradicionalmente como aquellos que
carecen de objeto directo y solo disponen de un argumento sujeto. Sin embargo,
semánticamente, la estructura de los intransitivos es más compleja y hoy en día
se reconocen dos subclases de intransitivos: inergativos e inacusativos.

\subsubsection{Verbos inergativos}\label{sec:07-inergativos}

Son aquellos en los que el sujeto desempeña el papel temático de agente. Se
agrupan en diferentes clases:

\begin{enumerate}[label=\alph*)]
  \item Verbos que seleccionan un sujeto encargado de realizar la acción que
        tiene lugar en el propio individuo: \forma{bromear, gesticular,
        rezongar, temblar}.
        \begin{exe}
          \ex Los encargados de la seguridad gesticularon con muecas divertidas.
        \end{exe}
  \item Verbos que implican la emisión de un sonido o sustancia: \forma{gritar,
        ladrar, llorar, sudar}.
        \begin{exe}
          \ex El cachorro ladró durante toda la noche.
        \end{exe}
  \item Verbos que señalan la manera de moverse, que pueden tener un uso
        claramente inergativo: \forma{bailar, caminar, correr}.
        \begin{exe}
          \ex Los egresados bailaron hasta el amanecer.
        \end{exe}
\end{enumerate}

\subsubsection{Verbos inacusativos}\label{sec:07-inacusativos}

El segundo grupo de verbos intransitivos es el de los inacusativos; en ellos, el
sujeto es tema o paciente. Los verbos inacusativos se agrupan en diversos tipos:

\begin{enumerate}[label=\alph*)]
  \item Verbos de causa interna, que nunca admiten la existencia de un agente o
        causa externa, ya que refieren a un proceso en el que el sujeto puede
        cambiar de manera innata: \forma{arder, madurar, estallar, florecer,
        germinar}.
        \begin{exe}
          \ex La bomba molotov no estalló.
          \ex Los cerezos florecieron durante la primavera.
        \end{exe}
  \item Verbos que denotan cambio de estado psíquico o físico y que poseen
        también usos transitivos: \forma{adelgazar, engordar, hervir, encoger,
        calentarse, congelarse}.
        \begin{exe}
          \ex Se congelaron las cañerías durante el invierno.
          \ex La ropa encogió después del lavado.
        \end{exe}
  \item Verbos de movimiento de dirección inherente cuando se emplean con
        sujetos \sig{inanimados}: \forma{acercarse, alejarse, partir,
        aproximarse}.
        \begin{exe}
          \ex El auto blindado se alejó rápidamente.
          \ex El avión partió con retraso.
        \end{exe}
  \item Otros verbos, como \forma{moverse, bambolearse, temblar}, también
        presentan usos inacusativos.
        \begin{exe}
          \ex La silla se bamboleaba.
          \ex El gatito temblaba por la anestesia.
        \end{exe}
  \item Verbos de existencia o aparición, que tienden a dar lugar a
        construcciones de inversión locativa, es decir, el complemento de lugar
        aparece en posición inicial, como se puede observar en los ejemplos
        \exref{ex:07-bosque} y \exref{ex:07-alegria}.
        \begin{exe}
          \ex\label{ex:07-bosque} En el bosque aparecieron fantasmas.
          \ex\label{ex:07-alegria} En esta casa falta alegría.
        \end{exe}
\end{enumerate}

\paragraph{Comportamiento sintáctico de los verbos inacusativos.} Las pruebas
para el reconocimiento de los verbos inacusativos se basan en la idea de que el
sujeto de las estructuras inacusativas posee las mismas características
semánticas del argumento objeto directo de un verbo transitivo, aunque
sintácticamente adquiere rasgos de concordancia con el verbo que lo hacen
funcionar como sujeto de la oración. Algunas de las pruebas que se han
presentado para reconocer esta clase de verbos son las siguientes:

\begin{itemize}
  \item Los sujetos de los inacusativos, al igual que los objetos de los
        transitivos, pueden aparecer en construcciones de participio absoluto,
        \exref{ex:07-barco} y \exref{ex:07-avion}, lo cual no es posible para
        los inergativos, \exref{ex:07-ladrado} y \exref{ex:07-caminado}.
\end{itemize}

\begin{exe}
  \ex\label{ex:07-barco} {}[Llegado el barco al muelle], subimos rápidamente.
  \ex\label{ex:07-avion} {}[Una vez aterrizado el avión], me tranquilicé.
  \ex[*]{\label{ex:07-ladrado} [Ladrado el perro], advertimos que algo pasaba.}
  \ex[*]{\label{ex:07-caminado} [Caminado Pedro], nos fuimos con los demás.}
\end{exe}

\begin{itemize}
  \item La prueba anterior no es posible con los verbos que implican una
        realización gradual, como \forma{aumentar, mejorar, crecer}, porque las
        construcciones de participio absoluto exigen verbos que expresen el
        estado final.\footnote{Los verbos que indican el estado final se conocen
        como télicos y se estudian en el \capref{8}.}
\end{itemize}

\begin{exe}
  \ex[*]{[Aumentada la nafta], disminuyó el tráfico.}
  \ex[*]{[Crecidas las plantas], empezamos a disfrutar del jardín.}
\end{exe}

\begin{itemize}
  \item Los verbos inacusativos nunca admiten objetos directos
        cognados\footnote{Se denominan objetos o complementos cognados aquellos
        que comparten la base léxica del verbo, como \forma{soñar un sueño
        apacible} o \forma{vivir una vida armoniosa}.} \exref{ex:07-llegada}, a
        diferencia de los inergativos, que sí los aceptan \exref{ex:07-carrera}.
\end{itemize}

\begin{exe}
  \ex[*]{\label{ex:07-llegada} Llegó [una llegada fácil].}
  \ex\label{ex:07-carrera} Corrió [una carrera asombrosa].
\end{exe}

\begin{itemize}
  \item Los participios de los verbos transitivos e inacusativos pueden
        funcionar como modificadores cercanos a los adjetivos: \forma{una
        promesa rota, una obra acabada, una persona desaparecida, un bebé recién
        nacido}, pero no sucede lo mismo con los participios de los inergativos:
        \juicio{*}{\forma{una mujer sonreída}}. \textcite{bello2002} llamó
        \termino{participios deponentes} a los de los verbos
        inacusativos.\pendiente{El original cita «Bello (1928)»}
  \item Los verbos inacusativos tienden a ser incompatibles con el sufijo
        nominalizador \forma{-dor}, puesto que tiene valor agentivo:
        \juicio{*}{\forma{desaparecedor}}, \juicio{*}{\forma{crecedor}},
        \juicio{*}{\forma{moridor}}. Los inergativos, en cambio, aceptan bien
        este sufijo: \forma{corredor, trabajador, caminador}.
\end{itemize}

\subsection{Verbos con alternancia causativa-anticausativa incoativa}\label{sec:07-alternancia}

Existen verbos que denotan un cambio de estado físico o psicológico, que poseen
usos inacusativos y provienen de verbos transitivos. Veamos algunos ejemplos:

\begin{exe}
  \ex
    \begin{xlist}
      \ex La inundación estropeó las cañerías.
      \ex\label{ex:07-estropearon} Las cañerías se estropearon.
    \end{xlist}
  \ex
    \begin{xlist}
      \ex El calor derritió los helados.
      \ex\label{ex:07-derritieron} Los helados se derritieron.
    \end{xlist}
\end{exe}

En los ejemplos \exref{ex:07-estropearon} y \exref{ex:07-derritieron}, el
sujeto tiene el papel temático de tema. Se trata de usos inacusativos en los que
el propio objeto desencadena la acción que señala el verbo, sin que intervenga
ningún agente \parencite{rodriguezramalle2005}. A los verbos que intervienen en
este paradigma se los suele llamar verbos de \termino{alternancia
causativa-anticausativa incoativa}. Se los denomina incoativos porque estas
construcciones inacusativas señalan el comienzo de un proceso, aunque suelen
terminar en un estado final o resultante. Esto sucede especialmente cuando
proyectan un clítico \forma{se}.\footnote{Al incorporar el clítico \forma{se},
estos verbos se convierten en pronominales cuasirreflejos. Estas construcciones
se estudian en el \capref{10}.}

\begin{recuadro}[Para ampliar]
En otras lenguas romances, como el francés y el italiano, donde estos verbos
fueron primeramente identificados, existen algunos diagnósticos para reconocer
los verbos inacusativos. En primer lugar, para formar los tiempos compuestos,
los verbos inacusativos seleccionan el auxiliar \forma{essere / être}
(\sig{ser}). En cambio, los inergativos seleccionan \forma{avere / avoir}
(\sig{haber}).

\begin{exe}
  \ex \gll Mara è partita alle dieci.\\
           Mara es partida a.las diez\\
      \glt \sig{Mara partió a las diez.}
  \ex \gll I bimbi sono ormai cresciuti.\\
           los niños son ya crecidos\\
      \glt \sig{Los niños han crecido ya.}
\end{exe}

Ahora bien, cabe señalar que en italiano todos los verbos inacusativos
seleccionan \forma{essere}, pero no todos los que seleccionan el auxiliar
\forma{essere} son inacusativos. Por ejemplo, los reflexivos se construyen con
\forma{essere}, pero no son inacusativos:

\begin{exe}
  \ex \gll Giovanni si era guardato allo specchio a lungo.\\
           Juan se era mirado en.el espejo a largo\\
      \glt \sig{Juan se miró en el espejo durante largo rato.}
\end{exe}

En segundo lugar, otra propiedad típica del objeto directo que el sujeto
pospuesto de los verbos inacusativos presenta en italiano es la posibilidad de
usar el pronombre partitivo \forma{ne}:

\begin{exe}
  \ex \gll Ne sono arrivati gli autori.\\
           de.ellos son llegados los autores\\
      \glt \sig{Llegaron los autores.}
\end{exe}

Pero este reemplazo no es posible con los verbos inergativos:

\begin{exe}
  \ex \gll {*}Ne hanno telefonato gli autori.\\
           de.ellos han telefoneado los autores\\
      \glt \sig{Llamaron por teléfono los autores.}
\end{exe}

En síntesis, podemos concluir que, en lenguas como el italiano y el francés, los
diagnósticos de inacusatividad tienen un carácter claramente sintáctico. Esto no
es tan evidente en el español actual. No obstante, sabemos que el español
medieval seleccionaba el auxiliar \forma{ser} para los verbos inacusativos:

\begin{exe}
  \ex Los moços somos idos a comer y nos han dexado solos. (Juan de Valdés,
      \textit{Diálogo de la lengua}, hacia 1535)\\
      \sig{Los mozos se han ido a comer y nos han dejado solos.}
\end{exe}
\end{recuadro}

\begin{ejercicio}\label{ej:07-inacusativos}
Considere el siguiente texto (A) y las oraciones de (B). Luego, según se muestra
en el ejemplo resuelto de abajo:

\begin{consignas}
  \item clasifique los verbos subrayados en transitivos, inergativos e
        inacusativos. Fundamente su respuesta con las pruebas adecuadas;
  \item analice las alternancias propuestas en (B) y explique qué cambios se
        producen en relación con los esquemas transitivos, intransitivos e
        inacusativos.
\end{consignas}

\textbf{A}

% Fragmento literario protegido: se incorpora a mano desde el original.
\textoaincorporar{los fragmentos de «Emma Zunz» que van de «De pronto,
alarmada, se levantó» hasta «\ldots{}la empezó a leer y la rompió» y de «El
hombre la condujo a una puerta» hasta «\ldots{}las partes que los forman», con
los verbos subrayados: \textit{levantarse, correr, abrir, dejar, ver, leer,
romper, conducir, haber, cerrarse, formar}}

\textbf{B}

\begin{exe}
  \ex
    \begin{xlist}
      \ex Abrió el cajón de la cómoda.
      \ex El cajón de la cómoda se abrió de repente.
    \end{xlist}
  \ex
    \begin{xlist}
      \ex El déficit empeoró la balanza comercial.
      \ex Las condiciones meteorológicas empeoraron en las últimas horas.
    \end{xlist}
  \ex
    \begin{xlist}
      \ex Una puerta se cerró.
      \ex Cerró la puerta con llave.
    \end{xlist}
\end{exe}

\begin{resuelto}
a) \forma{Levantarse} y \forma{correr} son verbos intransitivos. El sujeto de
ambos es agente y el verbo pertenece al grupo de los que seleccionan un sujeto
encargado de realizar la acción que tiene lugar en el propio individuo. b) La
oración 1a es transitiva: el sujeto es agente y posee un complemento objeto
directo tema. 1b es un caso de alternancia causativa inacusativa en la que el
verbo es inacusativo, pero proviene de un verbo transitivo, \forma{abrir}.
\end{resuelto}
\end{ejercicio}

\begin{recuadro}[Reflexión final]
En este capítulo nos hemos concentrado en la relación entre la semántica y la
sintaxis. Hemos definido los conceptos de argumento y predicado; luego, nos
hemos ocupado de la estructura temática y hemos definido las funciones
semánticas o papeles temáticos. A partir de allí, presentamos las funciones
sintácticas, específicamente aquellas incluidas en el sintagma verbal. Para
finalizar, desarrollamos los esquemas verbales transitivo e intransitivo y,
dentro de los últimos, hemos distinguido entre los inergativos y los
inacusativos.
\end{recuadro}

% --- Secciones finales -------------------------------------------------
\lecturasrecomendadas

\solucionario

\begin{solucion}{\ref{ej:07-agramatical}}
3. Es agramatical porque faltan dos argumentos, referidos a \sig{lo colocado} y
\sig{el lugar de colocación}. 4. Es agramatical porque falta un argumento
referido a \sig{lo que provocó el viento}. 5. Hay una incompatibilidad
léxico-semántica entre el verbo \forma{anunciar} y el argumento
\forma{gelatina}. 6. Es ambigua: a) es gramatical en la interpretación de
\sig{hizo un reclamo a los bancos}; b) es agramatical en la interpretación en
que falta el argumento referido a \sig{lo que reclamó el presidente}. 7. Es
agramatical porque falta un argumento referido al lugar. 8. Es agramatical
porque falta un argumento referido al lugar.\pendiente{El original numera estas
respuestas como (8) y (9)}
\end{solucion}

\begin{solucion}{\ref{ej:07-seleccion}}
a) 2 es gramatical, aunque hay un argumento implícito. 3 es gramatical, aunque
tiene dos argumentos implícitos. 5 es gramatical.\pendiente{Falta la preposición
\textit{a} ante \textit{su casa}: revisar el juicio} 6 es agramatical porque
\forma{a su aceituna} no constituye un argumento de lugar. b) \forma{Regresar}
[Arg.~1 SN, Arg.~2 SP].
\end{solucion}

\begin{solucion}{\ref{ej:07-valencia}}
a) Son bivalentes porque seleccionan dos argumentos: \forma{engañar, inquietar,
ingerir, firmar} y \forma{dirigir(se)}. El único monovalente, que selecciona un
solo argumento, es \forma{fallecer}. b) Los argumentos de \forma{engañar} son
\forma{el sello y el sobre} (sujeto) y el pronombre personal en caso acusativo
\forma{la} (COD). Los argumentos de \forma{inquietar} son \forma{la letra
desconocida} (sujeto) y el pronombre personal en caso acusativo \forma{la}
(COD). Los argumentos de \forma{ingerir} son \forma{el señor Maier} (sujeto) y
\forma{una fuerte dosis de veronal} (COD). El argumento de \forma{fallecer} es el
sujeto tácito, correferencial con el sujeto de \forma{ingerir}. Los argumentos
de \forma{firmar} son \forma{un compañero de pensión de su padre} (sujeto) y
\forma{la noticia} (COD). Los argumentos de \forma{dirigir(se)} son el sujeto
tácito, correferencial con el sujeto de \forma{firmar}, y \forma{la hija del
muerto} (COD).\pendiente{Revisar: \textit{a la hija del muerto} no es COD de
\textit{dirigirse}}
\end{solucion}

\begin{solucion}{\ref{ej:07-compadj}}
a) \forma{Preparó}: \forma{una sopa de tapioca} y \forma{unas legumbres}
(complementos). \forma{Comió}: \forma{temprano} (adjunto). \forma{Se acostó} y
\forma{se obligó a dormir}: no tienen complementos ni adjuntos. b) Los cuatro
verbos tienen el sujeto tácito o implícito; \forma{comió} tiene también el COD
implícito.
\end{solucion}

\begin{solucion}{\ref{ej:07-papeles}}
\forma{Una carta}: tema seleccionado por \forma{hallar}. \forma{El sello y el
sobre}: causa seleccionada por \forma{engañar}.\pendiente{El original dice «Una
carta y el sobre»} \forma{Su padre}: paciente seleccionado por \forma{morir}.
\forma{La letra desconocida}: causa seleccionada por \forma{inquietar}.
\forma{Emma}: agente seleccionado por \forma{leer}. \forma{El señor Maier}:
agente seleccionado por \forma{ingerir}. \forma{Un compañero de pensión de su
padre}: agente seleccionado por \forma{firmar}. \forma{La noticia}: tema
seleccionado por \forma{firmar}. \forma{A la hija del muerto}: meta seleccionada
por \forma{dirigir(se)}.
\end{solucion}

\begin{solucion}{\ref{ej:07-agente}}
\forma{Hilbert} no es agente porque rechaza la paráfrasis con \forma{hacer}
(\juicio{*}{\forma{Lo que hizo Hilbert fue ser\ldots}}) y la construcción final
(\juicio{*}{\forma{Hilbert era uno de los ponentes del Congreso para realizar su
investigación}}). \forma{La legendaria lista} no es agente porque es inanimada y
no admite la paráfrasis con \forma{hacer} (\juicio{*}{\forma{Lo que hizo la
legendaria lista\ldots}}) ni la construcción final (\juicio{*}{\forma{La
legendaria lista definió las matemáticas de la era moderna para lograr su
propósito}}).
\end{solucion}

\begin{solucion}{\ref{ej:07-mdo}}
Texto 1. \forma{Precisión matemática}: COD \sig{inanimado} y
\sig{específico}.\pendiente{Como COD escueto, \textit{precisión matemática} es
más bien inespecífico} Texto 2. \forma{Veraneos}: COD \sig{inanimado} e
\sig{inespecífico}. \forma{A su madre}: \forma{a} es marca de función (MDO); COD
\sig{animado} y \sig{específico}. \forma{La casita de Lanús}: COD
\sig{inanimado} y \sig{específico}. \forma{Los amarillos losanges de una
ventana}: COD \sig{inanimado} y \sig{específico}. \forma{El auto de prisión}:
COD \sig{inanimado} y \sig{específico}. \forma{El oprobio}: COD \sig{inanimado}
y \sig{específico}. \forma{Los anónimos}: COD \sig{inanimado} y
\sig{específico}.
\end{solucion}

\begin{solucion}{\ref{ej:07-seleccion2}}
b) \forma{Dar con} = \sig{encontrar}; \forma{optar por} = \sig{elegir};
\forma{conducir a}. \forma{Hombres del Nordstjärnan} está seleccionado por el
verbo \forma{dar (con)}. \forma{Por otro} está seleccionado por el verbo de
\sig{elección} \forma{optar}. \forma{A una puerta, a un turbio zaguán, a una
escalera tortuosa, a un vestíbulo, a un pasillo} y \forma{a una puerta que se
cerró} están seleccionados por el verbo de \sig{movimiento} \forma{conducir}.
\end{solucion}

\begin{solucion}{\ref{ej:07-inacusativos}}
a) \forma{Abrir, dejar, ver, leer, romper} y \forma{conducir} son verbos
transitivos. Todos seleccionan dos argumentos: un agente y un tema o paciente.
\forma{Haber} es un verbo impersonal transitivo que se construye con un COD con
el papel temático de tema. \forma{Cerrarse} es la variante inacusativa
pronominal del verbo transitivo \forma{cerrar}. \forma{Formar} es un verbo
transitivo con un COD pronominal.

b) 2. \forma{Empeorar}: alternancia transitiva / inacusativa sin incorporación
del clítico \forma{se}. En a) el sujeto es causa y en b) es tema. 3. Alternancia
transitiva / inacusativa \forma{cerrarse / cerrar}. El sujeto del verbo
inacusativo de a) se convierte en COD del verbo transitivo en b).
\end{solucion}
